% !TEX root = /home/eb/.vscode/extensions/qftrules.programme-colle/tmp/Exercice.tex%%%%%%%%%%%%%%%%%%%%%%%%%%%%%%%%%%%
\begin{exo}[PC][1][colle]{Tir à la carabine}
  Un homme tire à la carabine horizontalement une balle de $\SI{4,5}{mm}$ et de longueur $\SI{7}{mm}$ en plomb, de densité $d=\num{11}$.
  Estimer la portée maximale du tir.
\end{exo}
%%%%%%%%%%%%%%%%%%%%%%%%%%%%%%%%%%%

%%%%%%%%%%%%%%%%%%%%%%%%%%%%%%%%%%%
\begin{exo}[PC][3][oral]{Pendule de torsion}
  On considère une tige rigide $OA$, de longueur $\ell$, en liaison pivot parfaite en $O$ et à laquelle est fixée un masse $m$ en $A$. Le pendule évolue dans un plan vertical, et est relié à un ressort de torsion, de constante de torsion $C$, qui tend à ramener le pendule vers la position verticale. Le pendule est aussi soumis au poids.
  \begin{quest}
  \item Déterminer l'équation différentielle du mouvement de la masse $m$.
  \sol{
  TMC :
  $$
  m \ell^2 \ddot{\theta} = - C \theta + mg\ell \sin \theta \Rar \ddot{\theta} + \frac{C}{m\ell^2} \theta - \frac{g}{\ell} \sin \theta = 0.
  $$
  }
  \item Déterminer l'énergie potentielle du pendule pour toute position.
  \sol{
  Travail élémentaire associé au couple de torsion, $\delta W = - C \theta \d \theta = - \d E_p$ avec $E_p = \frac{1}{2}C \theta^2$. Pour la pesanteur, $\delta W = mg\ell \sin \theta \d \theta = - \d E_p$ avec $E_p = mg\ell(1-\cos \theta)$. L’énergie potentielle totale vaut
  $$
E_p = mg\ell(1-\cos \theta) + \frac{1}{2}C \theta^2.
  $$
  }
  \item Déterminer les positions d'équilibre, et leurs conditions d'existence, pour le pendule.
  \sol{
$\dd{E_p}{\theta} = C \theta-mg\ell\sin \theta = 0 $ pour $\theta = 0$ ou $\theta = \theta_0$ si $C<mg\ell$ (résolution graphique).
  }
  \item Déterminer si les positions d'équilibre précédentes sont stables ou instables (on pensera à séparer les différents cas).
  \sol{
  $\ddd{E_p}{\theta} = C-mg\ell\cos \theta$. Pour $\theta = 0$ , on trouve $\ddd{E_p}{\theta} = C-mg\ell$ stable ($>0$) si $C>mg\ell$ et instable ($<0$) si  $C<mg\ell$. C’est l’inverse pour l’autre position d’équilibre.
  }
\end{quest}
\end{exo}
%%%%%%%%%%%%%%%%%%%%%%%%%%%%%%%%%%%
